% ED26-BM-RAT2021-10-5: exact MIT 2021 chapter, printed 625--632.
\chapter{Structural Rationality: Acts, Practices and Forms of Life}
\label{sec:p11-rationality-10-5}

Nida-R\"umelin, Gutwald, and Zuber's theory of structural rationality
(TSR) evaluates an act in relation to patterns of agency that extend
beyond the agent's present preference satisfaction. Its subject is an
individual's act embedded in a course of life and in interactions with
other people. The word \emph{structural} therefore does not merely name
consistency among a person's simultaneous beliefs and intentions. That
different usage, represented here by Broome's
Definition~\ref{def:p11-reasons-structural}, is one of the chapter's
explicit contrasts \parencite[pp.~625 and 628--629]{NidaRumelinGutwaldZuber2021}.

The account combines an explanatory ambition with normative commitments.
It seeks to describe cooperation, commitments, regret, and the continuity
of persons more adequately than pointwise preference optimization. It
also takes some reasons connected with these practices to count in favor
of acting. Neither ambition proves the other. A description of what people
accept does not, without an argument about the practice and its reasons,
establish what they should accept.

\section{The particular structural account and its target}
\label{sec:p11-rat105-account}

\begin{hbdefinition}{def:p11-rat105-dimensions}{Three dimensions of structural rationality}
In this source's TSR, the \emph{diachronic dimension} concerns the place
of a present act in intended future aims, plans, and a continuing course
of action. The \emph{intrapersonal dimension} concerns the plurality of
reasons bearing on the person's action, rather than reducing all reasons
to present instrumental preference fulfillment. The \emph{interpersonal
dimension} concerns action as a contribution to practices and interactions
with other agents.

A \emph{structure} is an intended pattern of actions in which a singular
act has a place. A \emph{form of life} supplies the established practices
and contexts within which people give, assess, and accept reasons. The
authors' concluding rationality criterion is that the person's series of
acts forms an intended pattern that the person can accept as part of
their form of life. This is their practical-rationality proposal, not a
mathematical definition of a feasible set or a claim that every accepted
practice is justified.
\end{hbdefinition}

The summary identifies traditional rational choice theory (RCT) as
evaluating an action's suitability to satisfy desires at the moment of
choice, without adequately considering future plans or cooperation.
Section~1 presents a conflict between an individual-centered tradition
and one beginning with roles, practices, and institutions. It associates
classical and neoclassical economics with the first, and sociology,
cultural studies, and social psychology with the second. The authors
explicitly acknowledge exceptions, especially behavioral economics and
bounded-rationality approaches. This is an attributed description of
intellectual traditions, not a finding that every economist or every
social scientist follows one paradigm.

TSR's proposed reconciliation keeps the acting individual while
redescribing the individual act through the structures of a form of life.
It does not replace all persons with a collective agent. Nor does it
deny the practical use of instrumental reasoning: its claim is that
instrumental means--end assessment alone is too narrow to be an
all-embracing account of practical reason
\parencite[pp.~625--626]{NidaRumelinGutwaldZuber2021}.

The criticism needs a specified target. A decision model can encode
long-run plans or others' welfare in its consequences and preferences.
That mathematical possibility is recognized in the source's final
compatibility argument. The deeper dispute is whether the model explains
why a promise, a joint practice, or a continuing commitment counts as a
reason, and whether it has described the alternatives adequately. Merely
writing a larger utility function does not resolve that explanatory or
normative question.

\section{The historical and ethical motivation}
\label{sec:p11-rat105-ethical-core}

The authors interpret Aristotle's \emph{akrasia} as a singular act that
fails to fit an intended or desired course of action; they prefer lack
of self-control to the usual weakness-of-will gloss. In their account,
classical practical philosophy did not sharply separate morality and
rationality. A virtuous life realizes \emph{eudaimonia}, understood within
the political community rather than as an isolated individual's private
success. The good of the community supplies a context for assessing the
individual life. This provides an antecedent for structural thinking,
not a proof that ancient philosophy already supplied the present TSR
\parencite[p.~625]{NidaRumelinGutwaldZuber2021}.

Their Enlightenment narrative emphasizes the separation of rationality
and morality. In the Benthamite and later utilitarian line they discuss,
individual well-being or preference fulfillment supplies the basic value,
but moral obligation maximizes its total across persons rather than the
agent's own fulfillment. In RCT the object maximized need not be the
agent's well-being: it is the agent's actual subjective preference.
That distinction prevents selfish welfare from being read into every
utility representation.

In their Kantian comparison, pragmatic imperatives concern one's
well-being, whereas moral motivation is respect for a law requiring a
maxim suitable for universal adoption. Universal maxims already locate
an act in a structure. The source raises the formalism objection:
appropriate redescription can leave the universal-law test too
underdetermined to settle particular choices. It connects this criticism
with Hegel and communitarian discussions. Against total-utility
maximization it invokes Rawls's separateness of persons: summing
fulfillment can fail to preserve individual motives and life plans.
These are the authors' critical interpretations and objections, not
independently established verdicts on the whole Kantian or utilitarian
tradition \parencite[pp.~625--626]{NidaRumelinGutwaldZuber2021}.

The chapter further argues that ordinary game theory already introduces
one structural element, since equilibrium relates each decision to the
other decisions. Harsanyi's surveyed rule-utilitarian or ethical-Bayesian
account makes the moral criterion coordinated practice optimizing a
common value function, rather than separately maximizing total
fulfillment at each act. Its specified common value is the sum of
preference fulfillment. That is a particular moral aggregation premise;
it is neither the definition of all cooperation nor a consequence of
having utility numbers. A rule's suitability for joint adoption and an
act's unilateral effects ask different questions
\parencite[p.~626]{NidaRumelinGutwaldZuber2021}.

\section{The prisoner's dilemma and the description of cooperation}
\label{sec:p11-rat105-cooperation}

The source's transition from decision theory to game theory starts with
values of consequences and subjective probabilities over circumstances.
In interaction, those circumstances include the other participants'
decisions, possibly together with chance events. Treating them as part
of one's environment retains an individual optimization problem; asking
how the choices fit together adds a structural question. The source
describes noncooperative choices as mutually independent. This must not
be read as the definition of stochastic independence: the independent
draws below are an explicit mixed-strategy modeling assumption, whereas
the players' choices can depend strategically on expectations about one
another \parencite[p.~626]{NidaRumelinGutwaldZuber2021}.

\begin{hbexample}{def:p11-rat105-source-pd}{The source's exact cooperation problem}
Players \(A\) and \(B\) choose cooperation \(C\) or defection \(D\).
The source's Table~10.5.1 supplies these payoffs, with the row player's
value first:
\[
 \begin{array}{c|cc}
       &C&D\\ \hline
 C&(3,3)&(1,4)\\
 D&(4,1)&(2,2).
 \end{array}
\]
For the following formal calculation these are the complete cardinal
utilities to be maximized, players compare expected utilities, and
mixed draws are independent as in
Definition~\ref{def:p17-finite-strategic-game}. If instead these numbers
represent only money, they need not exhaust either player's values.
\end{hbexample}

\begin{hbtheorem}{thm:p11-rat105-pd}{proved=annex:p11.rat105.pd}{Optimistic beliefs do not remove the original dominance}
In the stipulated game, \(D\) strictly dominates \(C\) for both players
against every opponent mixture. Its unique independent mixed Nash
equilibrium is \((D,D)\). Mutual cooperation nevertheless improves both
utilities, from \((2,2)\) to \((3,3)\). If \(q\in[0,1]\) is a
player's belief that the opponent chooses \(C\), the values of \(C\)
and \(D\) are respectively \(1+2q\) and \(2+2q\). Even certainty
of the other's cooperation leaves a utility gain of one from defection.
\end{hbtheorem}

The mathematical result distinguishes unilateral improvement with the
other's choice fixed from comparison of joint outcomes. It does not make
efficiency, consensus, and rationality identical. Neither participant can
select the mutually cooperative profile alone. Nash equilibrium is the
no-profitable-unilateral-deviation condition of
Definition~\ref{def:p17-nash-equilibrium}; it is not, without further
epistemic premises, the set of all choices rational relative to arbitrary
beliefs. The source's general equation of game-theoretic rational
decisions with equilibrium therefore needs this qualification. In its
particular dilemma, dominance already forces defection under the stated
utility criterion, so no general epistemic-equilibrium identification
is needed \parencite[p.~626]{NidaRumelinGutwaldZuber2021}.

The authors' actual cooperative speaker wants the two players both to
cooperate and expects the other to cooperate. They argue that a general
theory neutral about motivation should make room for this motive.
The speaker conceives their own choice as doing their part in a joint
action, hoping the other contributes their part. The description of an
option as cooperative or defective refers to this structure, not merely
to a probability attached to a final payoff. TSR's requirement of a
comprehensive description of strategic options is meant to preserve that
relationship \parencite[p.~627]{NidaRumelinGutwaldZuber2021}.

The example exposes a real fork. If the four numbers are the player's
complete utilities and only unilateral expected-utility improvement
counts, the quoted belief does not justify \(C\): even at \(q=1\),
three is less than four. TSR must change the relevant practical
description, the standard of assessment, or the preference representation.
It cannot obtain cooperation by recalculating the unchanged game more
optimistically. If the numbers are merely material rewards, an additional
preference for participation in mutual cooperation may change the
effective game. Whether such a preference is present is descriptive;
whether the joint practice supplies a good reason is normative.

Shared or collective intention is one possible reconstruction. The
source surveys Hurley, Tuomela, Gilbert, Bratman, and Regan, and describes
Tuomela's group ethos as changing the individual utilities. These accounts
are not interchangeable definitions of one mechanism: the chapter does
not construct or prove all of them. Its asserted interpretation of
cooperative behavior is that agents favor \(C\) as a contribution toward
\(\langle C,C\rangle\), which also makes disappointment when the other
defects intelligible. Disappointment alone cannot establish that this
was the actual motive or the uniquely best explanation.

\subsection{What the cited experimental evidence can establish}

The chapter reports substantial cooperation, including in one-shot
tasks, and recommends interpreting it as structurally rational. Its
citation group is heterogeneous: finitely repeated and sequential
dilemmas, experimental papers, a theoretical incomplete-information
model, and broader works on cooperation do not all constitute evidence
for the identical one-shot task. The claim that TSR explains these
findings best is the authors' philosophical assessment. The selected
chapter supplies no comparative model fit, operational TSR measure, or
population-level estimate that establishes it
\parencite[pp.~626--627]{NidaRumelinGutwaldZuber2021}.

One of its cited primary experiments permits a bounded check.
Barreda-Tarrazona and colleagues recruited undergraduate volunteers at
Universitat Jaume I in Spain in December 2014. Of 178 first-session
participants, 162 returned for the prisoner's-dilemma sessions. They
classified participants using giving in a dictator game and scores on
a 40-item, 20-minute abstract-reasoning test, creating four groups around
the median scores. Participants were randomly paired within their group,
without being told of the grouping. Each person first played ten
anonymous stranger-paired one-shot encounters, then three ten-round
games with a continuing partner within each game. Thus \emph{one-shot}
here means no continuing relationship with a known partner, not no
previous laboratory encounters
\parencite[pp.~3--5]{BarredaTarrazonaEtAl2017Cooperation}.

Their monetary payoff pairs were \((20,20),(0,28),(28,0),(10,10)\)
in the one-shot task, with all stage rewards divided by ten in repeated
tasks. Incentivized forecasts of partner cooperation and of session-wide
cooperation were elicited. The paper distinguishes an individual's
cooperation from successful cooperation by both members of a pair.
It reports declining cooperation across the one-shot encounters and
more sustained cooperation, followed by an end-of-game drop, in the
repeated tasks. Almost no cooperation occurred when a participant
expected defection: the reported one-shot unconditional-cooperation
rates average 1.5\% of decisions for low-giving participants and 2.8\%
for high-giving participants
\parencite[pp.~8--12]{BarredaTarrazonaEtAl2017Cooperation}.

The uncertainty and scope matter. Its Table~6 random-effects logit
model reports a one-shot reasoning-score coefficient of \(-0.09\)
(standard error \(0.03\)) and a dictator-giving coefficient of
\(0.20\) (standard error \(0.09\)); the first repeated game's
reasoning coefficient is \(0.06\) (standard error \(0.03\)).
These are coefficients on log odds, not percentage-point effects.
The 1,620 observations per task are ten decisions by each of 162
people, not 1,620 independently recruited participants. The paper
reports significance thresholds and standard errors rather than
confidence intervals for these displayed coefficients. Its groupings
measure traits rather than randomly changing intelligence or altruism,
and its task order is fixed. Trait associations and the contrast between
later repeated games and earlier one-shot encounters therefore do not
isolate causal effects of the traits or of repetition alone.

The authors discuss altruism, expectations, reputation building, and
adaptation to the game as explanations. Their experiment does not
measure TSR's intended life structures or rule out these alternatives.
It is evidence of cooperation under a specified task and sample, not
a validation of the normative superiority of cooperation, the universal
rationality of \(C\), or a generalization to every social practice.
No raw-data reanalysis or independent replication is supplied here.
The wider citation list in the 2021 chapter is not silently promoted to
a homogeneous empirical result
\parencite[pp.~9--12, especially Table~6]{BarredaTarrazonaEtAl2017Cooperation}.

\section{Meta-preferences and continuing personhood}
\label{sec:p11-rat105-metapreferences}

The source considers meta-preferences as a way to introduce structural
evaluation without abandoning the belief--desire model. In its
reconstruction of Sen, one level retains the egoist dilemma ordering;
an assurance-game ordering favors cooperation when the other cooperates
but resists exploitation by a defector; an other-regarding ordering
favors cooperation regardless of the other's choice. Conditional and
unconditional cooperation are thus different rankings, not two solutions
to the unchanged dominant-defection problem.

For a worked contrast, compare the original game with an assurance
ordering: a player prefers \((C,C)\) to \((D,C)\), but prefers
\((D,D)\) to \((C,D)\). Now expectations about the other's action
can change the preferred own action. In the original matrix, defection
is preferred in both cases. An unconditional cooperator instead ranks
\((C,C)\) over \((D,C)\) and \((C,D)\) over \((D,D)\).
These comparisons describe the source's three types without pretending
that it provides cardinal values for each meta-ordering. They do not
prove which ordering is morally justified
\parencite[p.~627]{NidaRumelinGutwaldZuber2021}.

Frankfurt's surveyed version concerns second-order volitions about which
first-order action-guiding desires should constitute one's will. The
person can distance themselves from a desire and care about what guides
them. The contrasted wanton follows first-order desires without this
concern for their will. TSR uses these structural features to explain
continuity through time and interaction: persons evaluate their
action-guiding preferences rather than merely optimizing anew at every
moment.

The authors argue that an agent who finds their preferences structurally
incoherent can wish for different first-order preferences. Such evaluation
may change them; it need not actually do so. If it does not, the
structural judgment and pointwise rational choice can conflict. Adding
the word \emph{meta} does not guarantee a consistent hierarchy or
explain how a higher-order judgment prevails. The source's discussion
supplies a conceptual account of self-evaluation, not a population study
of second-order volitions or a complete dynamic revision algorithm
\parencite[p.~627]{NidaRumelinGutwaldZuber2021}.

\section{Reasons, promises, roles and individual rights}
\label{sec:p11-rat105-reasons}

The authors argue that the sharp rationality/morality division collapses
when rational action is assessed through the plurality of practical
reasons. The bridge requires the premise that moral and other
nonconsequentialist considerations can be genuine practical reasons.
They would then belong within practical assessment rather than being
external restrictions on an independently exhaustive rationality
criterion. The following alternatives expose ways to reject, limit,
or implement that premise
\parencite[p.~628]{NidaRumelinGutwaldZuber2021}.

\begin{hbdefinition}{def:p11-rat105-four-responses}{Four responses to nonconsequentialist reasons}
The source distinguishes four reactions to reasons that do not simply
maximize consequences. A \emph{consequentialist} response excludes such
prima facie considerations as nongenuine or incapable of integration.
A \emph{narrow-rationality} response limits RCT's jurisdiction and leaves
moral or other good reasons outside it. A \emph{revisionist} response
changes constitutive RCT axioms, as in the surveyed resolute-choice
approach of McClennen. TSR's \emph{redesign} response reinterprets basic
rational-choice concepts to include nonconsequentialist reasons.
\end{hbdefinition}

The first response risks excluding central kinds of reasons; the second
concedes that RCT is not a complete theory of practical rationality.
The third alters the formal framework. The fourth tries to preserve a
mathematical core while changing its practical interpretation. These
are different strategies, not conclusions that follow from the bare
existence of several kinds of considerations. The fourth owes an
adequate description of those reasons and an argument for their standing
\parencite[p.~628]{NidaRumelinGutwaldZuber2021}.

\begin{hbexample}{def:p11-rat105-promise}{A promise within a continuing practice}
The source's promisor has a reason to keep a promise even when keeping
it produces less-than-optimal consequences for the promisor or the
promisee. The relevant premise is that giving a valid promise within the
institution supplies a duty or practical reason. The authors invoke
Searle's speech-act analysis for the institution's normative elements.

TSR describes keeping it as sustaining diachronic consistency in an
interpersonal practice: the undertaking allows other people to relate
their plans to the promisor's anticipated acts. The example therefore
includes the original commitment, its continuing significance, and the
others whose plans depend on it. Merely comparing the immediate
outcomes misses the reason the act is described as promise keeping.
\end{hbexample}

This is not a proof that every promise overrides every consequence.
Validity, the content promised, conflicting obligations, and possible
release from the promise can affect the reason. A promise to perform a
wrongful act is a limit case requiring justification, not automatically
a requirement because it forms a stable pattern. Predictability can
explain a practice's effects without supplying the whole reason for its
normative authority.

The authors apply the same structural interpretation to duties connected
with being a parent, teacher, or politician: such roles have temporally
and socially connected responsibilities. Individual rights, in contrast
to simply maximizing a common sum, secure authorship of one's acts in
social contexts and boundaries within which the person can act. The
account must preserve both social connections and the individual whose
agency is protected. It cannot infer that any established role deserves
obedience merely because it coordinates behavior.

Their conclusion is that structure is not an appraisal function imposed
from outside the practice of giving and accepting reasons. It
systematizes that practice. Their stronger assertion that consistent
practice, personhood, and collective rationality require reference to
structures is a philosophical claim supported here by the roles of
continuity, commitment, and coordination. The chapter does not establish
an impossibility theorem for every rival theory
\parencite[p.~628]{NidaRumelinGutwaldZuber2021}.

\section{The instrumental and attitude-coherence contrast}
\label{sec:p11-rat105-instrumental}

The source describes a common analytic use of structural rationality as
a relational consistency requirement among attitudes, often assessed at
one time, without substantive criticism of their contents. Its contrast
is not that relations cease to matter in TSR. It is that TSR additionally
includes interpersonal structures and substantive reasons. The authors
classify the narrower usage as instrumental; that classification is
their interpretation, not a theorem that every attitudinal-coherence
account accepts only instrumental reasons
\parencite[pp.~628--629]{NidaRumelinGutwaldZuber2021}.

Bratman is an important exception to a simple individual/social split.
The source's account of his modest sociality extends individual agency's
conceptual and normative resources to shared intentions and planning.
Friendship, love, conversation, dancing together, and marriage illustrate
practices that need social description. Merely intending the same thing
is not enough for shared action. The participants need coordinating
plans and subplans, with the appropriate common knowledge; cooperation
can fail if a required contribution is absent. Intentions function as
plan states organizing conduct across time and across participants.
This overlaps TSR's concerns without supplying TSR's full substantive
account of reasons.

\begin{hbexample}{def:p11-rat105-intention-conflict}{The source's family--manuscript conflict}
The Langlois example has an agent intending both to spend the weekend
with their family and to finish a new manuscript by Sunday, while
believing the two plans cannot both be carried out. The irrational
combination, on the surveyed coherence account, lies in the relation
between these intentions and that belief. Neither goal is criticized
merely for its content.

Removing the belief in incompatibility, changing one plan's timing, or
abandoning one intention changes the example. It is therefore not
evidence that having two demanding projects is always irrational.
TSR can ask the further substantive questions about promises, family
roles, and a continuing life plan; those questions are not settled by
the relational inconsistency alone.
\end{hbexample}

The source credits Scanlon's explicit structural-irrationality usage
while note~1 records earlier German uses in the critique of
consequentialism and later developments by Nida-R\"umelin. It also
surveys belief inconsistency, enkratic requirements, and Broome's 1999
analysis of arrangements of attitudes. Its note~2 points to Kiesewetter's
2017 analysis. These provenance and continuation notes do not erase the
different meaning selected for the present chapter.

Broome's local treatment keeps rationality's mental dependence distinct
from correct response to actual normative reasons and argues against
identifying the two. TSR explicitly uses reasons embedded in practices
and seeks to dissolve the rationality/morality separation at that level.
The two treatments therefore differ on what is assessed and what can
constitute a reason for the assessment. A comparison must not make
Broome's mental-supervenience premise an unstated TSR axiom or make
TSR's practical-reasons criterion Broome's definition
\parencite[pp.~628--630]{NidaRumelinGutwaldZuber2021}.

\section{Compatibility with expected utility: what is and is not proved}
\label{sec:p11-rat105-ramsey}

Section~7 argues that utility representation is neutral about motivation.
Qualitative preferences satisfying the appropriate axioms can admit
quantitative representations even when their motives are altruistic,
deontological, or connected with social practices. The authors name
Ramsey, von Neumann and Morgenstern, Marschak, and Luce and Raiffa in
the historical development, and mention consistency, transitivity,
monotonicity, and lottery continuity. The selected chapter does not
state one complete axiom system, one exact domain, or its representation
proof. Different representation theorems must not be conflated by
calling all of them the utility theorem
\parencite[pp.~629--630]{NidaRumelinGutwaldZuber2021}.

For an exact supplied interface, the finite-prize model of
Theorem~\ref{thm:p17-finite-vnm} uses preferences on \emph{all}
lotteries over a finite nonempty consequence set, together with
completeness, transitivity, mixture independence, and its specified
continuity condition. It establishes the representation and the full
proof is at its reciprocal annex item. None of those axioms requires
that the consequence's value be selfish pleasure or money. Conversely,
having a virtuous motive does not establish that one's whole lottery
preference satisfies independence or continuity. Applying a
representation theorem remains conditional on its actual hypotheses.

\begin{hbtheorem}{thm:p11-rat105-value}{proved=annex:p11.rat105.value}{Structural values can have an expected-value representation}
Let \(Z\) be any finite nonempty set of fully described consequences
and let \(v:Z\to\mathbb R\) be a supplied value function. For a
lottery \(p\) on \(Z\), put
\[
 V_v(p)=\sum_{z\in Z}p(z)v(z).
\]
Comparing lotteries by \(V_v(p)\geq V_v(r)\) yields a complete,
transitive preference satisfying mixture independence and the continuity
condition of Definition~\ref{def:p17-lottery-continuity}. For real
\(a>0\) and \(b\), replacing \(v\) by \(av+b\) changes
\(V_v(p)\) to \(aV_v(p)+b\) and preserves all comparisons.

In particular, take the four joint action profiles of
Example~\ref{def:p11-rat105-source-pd} as \(Z\) and define an
editorial structural value by
\(v(\langle C,C\rangle)=1\) and \(v(z)=0\) at the other three
profiles. Against an opponent who cooperates with probability \(q\),
the expected structural values of one's \(C\) and \(D\) are \(q\)
and zero. Cooperation is strictly preferred when \(q>0\); at
\(q=0\) both are indifferent.
\end{hbtheorem}

This explicit model proves a limited compatibility possibility: a
preference for contributing to mutual cooperation need not violate the
expected-value mathematics. It is an editorial example, not an elicited
preference of the source's experimental participants, nor a replacement
for the original game's utilities. It assigns value to the jointly
cooperative profile and consequently changes the strategic problem.
The source's note~3 about positive changes of origin and scale is the
same affine identity; full converse uniqueness on the lottery domain is
Theorem~\ref{thm:p17-vnm-affine-uniqueness}. An arbitrary increasing
change of ordinal values need not preserve expected-value comparisons.

The source's more general conditional claim requires both coherent
probabilities and preferences satisfying the relevant Ramsey
representation conditions. Kolmogorov's probability axioms alone do
not provide those preference conditions, do not identify the agent's
beliefs with the relevant probabilities, and do not say which reasons
are good. A representation gives numerical form to comparisons;
selecting a maximizing alternative additionally uses those comparisons
as the decision standard. The source describes the retained core as
coherence between epistemic attitudes (probabilities) and prohairetic
attitudes (preferences). Calling the result subjective value rather
than utility avoids a hedonistic suggestion, but does not establish a
normative justification by a change of words.

The chapter therefore defends compatibility with the mathematical core
of RCT while rejecting its claimed monopoly on practical reasons. It
does not show that all structurally motivated choices maximize the old
unrevised utility, that every form of life has a unique utility, or that
any expected-value maximizer thereby satisfies TSR.

\section{Description, normativity and use}
\label{sec:p11-rat105-use}

The source's final section treats both RCT and TSR as descriptive and
normative: economic rationality is used to explain conduct and to advise
it; TSR begins with observable practice while criticizing structurally
irrational conduct. Regret and repentance can reveal the person's
evaluation of acts against a continuing pattern. Duties may arise in
roles, personal commitments, prescriptions, and rights. The authors
oppose reducing all of them to optimizing states of affairs and insist
that theory should not dictate from outside which reasons constitute a
form of life. Their note~4 connects the autonomy issue with Korsgaard
\parencite[p.~630 and nn.~1--4]{NidaRumelinGutwaldZuber2021}.

The justification is immanent: participants start with their practices
of giving and accepting reasons, and can argue that a particular reason
is inadequate or conflicts with a more fundamental one. The authors
deny that participants can step entirely outside the normative frame of
their form of life. This is a philosophical position about practical
reason, not an experimental result. It permits criticism within the
practice; it does not follow that popularity validates a reason. An
oppressive but stable role, incompatible commitments, or disagreement
about what counts as a valid promise tests the theory's capacity to
explain revision and substantive criticism. The selected chapter supplies
no universal resolution rule for such conflicts.

\begin{hbmethod}{def:p11-rat105-structure-check}{Using the source's account without changing the question}
Describe the singular act, the intended course of action, the other
participants and contributions, and the practice in which reasons are
offered. State separately the temporal plan, the person's competing
reasons, and the interpersonal structure. Check whether the proposed
act actually fits that structure rather than merely wishing for a
collective outcome.

Distinguish a content-based practical reason from a consistency relation
among attitudes. In the family--manuscript case, identify the two
intentions and the incompatibility belief; in the promise case, identify
the undertaking and the claimed institutional reason. Do not let either
example substitute for the other.

If applying a choice model, say whether its payoffs are money or full
utilities, whether it ranks own actions or joint structures, and which
probability and preference hypotheses apply. Test the unchanged
dilemma: optimistic beliefs alone cannot make its dominated action
maximize the stated utility. If a new value is introduced, show the
changed rankings and explain the reason for adopting it.

For a descriptive claim, specify the sample, task, operational measure,
design, uncertainty, and alternatives as required by the evidence record
in Method~\ref{def:p19-evidence-record}. For a normative conclusion,
state why the commitment, role, right, or practice supplies a reason;
observed cooperation, consensus, or efficient outcomes do not supply
that premise by themselves.
\end{hbmethod}

This procedure is an editorial application of the chapter's distinctions,
not its claimed universal decision algorithm. Its success case is a
valid promise whose fulfillment coheres with an accepted continuing
practice, with the actual normative premise and possible conflicts
made explicit. Its failure cases include a fulfilled but wrongful
promise, a plan whose necessary partner never contributes, and calling
\(C\) an expected-utility best response while retaining the original
matrix. These diagnose different failures rather than collapsing them
under a single rationality score.

The comparison with Schmid's
Chapter~\ref{sec:p17-rationality-10-2} is especially close but not an
identity. Schmid distinguishes effectiveness or efficiency of participant
cooperation from consistency of a corporate agent's attitudes; team
reasoning adopts a shared point of view and identifies one's part in a
joint choice. TSR instead organizes singular acts through diachronic,
intrapersonal, and interpersonal structures of practical reasons. A
shared intention is one possible reconstruction in its cooperation
discussion, not a requirement that every rights-holder or promisor be
part of one corporate mind. Neither account may identify consistency
with good objectives or joint success with moral correctness.

Broome's Chapter~\ref{sec:p11-reasons-and-rationality} preserves a
different mental and attitudinal evaluand and disputes an actual-reasons
identity. The finite P17 theories prove formal consequences of stated
preference and strategic assumptions. P19's
Definition~\ref{def:p19-benchmark} separates a standard from observed
conformity. Together these interfaces locate the present account without
making one source's normative criterion a universal definition of
rationality. The mathematical consequences supplied here have full
annex proofs; the source's practical-reasons, historical, and explanatory
claims retain their argued and attributed philosophical standing.

\index{structural rationality!Nida-R\"umelin, Gutwald and Zuber}
\index{rationality!diachronic}
\index{rationality!interpersonal}
\index{meta-preferences}
\index{promise keeping!structural reasons}
\index{prisoner's dilemma!structural rationality}
